\begin{question}

  Construct $\ket{\vectorbold{S}\cdot\vectorbold{\hat{n}}; +}$ such that
  \begin{equation*}
    \vectorbold{S}\cdot\vectorbold{\hat{n}} \ket{\vectorbold{S}\cdot\vectorbold{\hat{n}}; +}
    = \left(\frac{\hbar}{2}\right)\ket{\vectorbold{S}\cdot\vectorbold{\hat{n}}; +}
  \end{equation*}
  where $\vectorbold{\hat{n}}$ is characterized by the angles shown in
  the figure. Express your answer as a linear combination of $\ket +$
  and $\ket -$

  [\emph{Note:} The answer is
    \begin{equation*}
      \cos(\frac \beta 2)\ket + + \sin(\frac \beta 2) e^{i\alpha}\ket -
    \end{equation*}
    But do not just verify that this answer satisfies the above
    eigenvalue equation. Rather, treat the problem as a
    straightforward eigenvalue problem. Also do not use rotation
    operators, which we will introduce later in this book.]
  
\end{question}

\begin{answer}

  Grumble grumble grumble...stupid problem...grumble grumble...already
  saw rotation operators...baa humbug...grumble...

  sigh...ok...yes...this problem is really annoying...but I should
  probably put in some real annotations anyway.  My buddy Eric says
  that it's going to be a useful derivation later.  Still...at times
  like this, I really really like the eig function in GNU Octave.

  It's worth noting that what is happening is that we're constructing
  an $S_n$ operator where $\hat{n}$ is along any arbitrary 3D
  direction.  Putting it another way, what happens if you give your
  Stern Gerlach device to a drunk baby?\footnote{I said I was going to
  put in real annotations, I didn't say that I wasn't going to be
  surly about it.}  Even though it's a random orientation, you're
  still going to need a few things to be true, such as the eigenvalues
  needing to still be $\frac \hbar 2$\footnote{SPOILERS!!!}.

  The notation that Sakurai uses is a little heavyweight, so let's
  lighten it up a bit.

  \begin{equation}
    A \equiv \vectorbold{S}\cdot\vectorbold{\hat{n}}
  \end{equation}

  \begin{equation}
    \ket{\alpha} \equiv \ket{\vectorbold{S}\cdot\vectorbold{\hat{n}}; +}
  \end{equation}

  They give us one of the Eigenvalues right at the beginning, but
  what's the fun in that.  If we're going to embrace the painful
  irritation of computing everything by hand then let's do it!
  
  \begin{equation}
    \lambda_+ \equiv \left(\frac \hbar 2\right)
  \end{equation}

  Finally, we arrive at our characteristic equation like we all
  remember from childhood...ahhh those ``happy'' days...
  
  \begin{equation}
    A\ket{\alpha} = \lambda\ket{\alpha}
  \end{equation}

  Referencing the figure in the book\footnote{It's a vector graphic in
  the PDF and I didn't feel like either screenshotting it or
  recreating it, so just look it up in your book...you...\emph{did}
  buy the book right? RIGHT?}, we can construct our $\hat{n}$ which is
  expressed in spherical coordinates
  
  \begin{equation}
    \vectorbold{\hat{n}}
    \equiv
    \begin{bmatrix}
      \cos{\alpha}\sin{\beta}\\
      \sin{\alpha}\sin{\beta}\\
      \cos{\beta}\\
    \end{bmatrix}
  \end{equation}

  We recall some basic definitions
  
  \begin{equation}
    \tag{1.111a}
    S_x = \frac \hbar 2 (\ket - \bra + + \ket + \bra -)
  \end{equation}

  \begin{equation}
    \tag{1.111b}
    S_y = i \frac \hbar 2 (\ket - \bra + - \ket + \bra -)
  \end{equation}
  
  \begin{equation}
    \tag{1.90}
    S_z = \frac \hbar 2 (\Lambda_+ - \Lambda_-) = (\ket + \bra + - \ket - \bra -)
  \end{equation}

  And now we take the giant linear combination
  
  \begin{equation}
    \begin{split}
      A
      &\equiv \vectorbold{S}\cdot\vectorbold{\hat{n}}\\
      &= \cos{\alpha}\sin{\beta}S_x + \sin{\alpha}\sin{\beta}S_y + \cos{\beta}S_z\\
      &= \left(\frac \hbar 2\right)\left\{
        \cos{\alpha}\sin{\beta}(\ket - \bra + + \ket + \bra -)
        + i\sin{\alpha}\sin{\beta}(\ket - \bra + - \ket + \bra -)
        + \cos{\beta}(\Lambda_+ - \Lambda_-)
        \right\}\\
      &= \left(\frac \hbar 2\right)\left\{
        \sin{\beta}[\cos{\alpha} + i\sin{\alpha}](\ket - \bra +)
        + \sin{\beta}[\cos{\alpha} - i\sin{\alpha}](\ket + \bra -)
        + \cos{\beta}\Lambda_+
        - \cos{\beta}\Lambda_-
      \right\}\\
      &= \left(\frac \hbar 2\right)\left\{
        \sin{\beta}e^{i\alpha}(\ket - \bra +)
        + \sin{\beta}e^{-i\alpha}(\ket + \bra -)
        + \cos{\beta}\Lambda_+
        - \cos{\beta}\Lambda_-
      \right\}\\
    \end{split}
  \end{equation}

  Now that we have our 4 basic components ($++, +-, -+, --$), we
  recall how the book says we're supposed to cram them together.
  
  \begin{equation}
    \tag{1.73}
    X \dot{=}
    \begin{pmatrix}
      \mel{\alpha^{(1)}}{X}{\alpha^{(1)}} & \mel{\alpha^{(1)}}{X}{\alpha^{(2)}} & \ldots \\
      \mel{\alpha^{(2)}}{X}{\alpha^{(1)}} & \mel{\alpha^{(2)}}{X}{\alpha^{(2)}} & \ldots \\
      \vdots & \vdots & \ddots\\
    \end{pmatrix}
  \end{equation}

  We do some cramming and come back with our representative matrix
  that we can use to solve the Eigenvalue problem.
  
  \begin{equation}
    \inlabel{A-finish}
    \begin{split}
      A
      &\equiv \vectorbold{S}\cdot\vectorbold{\hat{n}}\\
      &= \left(\frac \hbar 2\right)\left\{
        \sin{\beta}e^{i\alpha}(\ket - \bra +)
        + \sin{\beta}e^{-i\alpha}(\ket + \bra -)
        + \cos{\beta}\Lambda_+
        - \cos{\beta}\Lambda_-
      \right\}\\
      &\dot{=}
      \begin{bmatrix}
        \mel{+}{A}{+} & \mel{+}{A}{-}\\
        \mel{-}{A}{+} & \mel{-}{A}{-}\\
      \end{bmatrix}\\
      &=
      \left(\frac{\hbar}{2}\right)
      \begin{bmatrix}
        \cos{\beta} & \sin{\beta}e^{-i\alpha}\\
        \sin{\beta}e^{i\alpha} & -\cos{\beta}\\
      \end{bmatrix}\\
    \end{split}
  \end{equation}

  Fleshing out our characteristic equation, we find that
  
  \begin{equation}
    \begin{split}
      \vectorbold{0}
      &= (A - \lambda I)\ket{A+}\\
      &\dot{=}
      \left\{
      \left(\frac{\hbar}{2}\right)      
      \begin{bmatrix}
        \cos{\beta} & \sin{\beta}e^{-i\alpha}\\
        \sin{\beta}e^{i\alpha} & -\cos{\beta}\\
      \end{bmatrix}
      - 
      \begin{bmatrix}
        \lambda & 0\\
        0 & \lambda\\
      \end{bmatrix}
      \right\}
      \ket{A+}\\
      &=
      \left(\frac{\hbar}{2}\right)
      \begin{bmatrix}
        \cos{\beta} - \frac{2}{\hbar}\lambda & \sin{\beta}e^{-i\alpha}\\
        \sin{\beta}e^{i\alpha} & -\cos{\beta} - \frac{2}{\hbar}\lambda\\
      \end{bmatrix}
      \begin{bmatrix}
        \alpha_+\\
        \alpha_-\\
      \end{bmatrix}
      \\
      &\Rightarrow\\
      &0 =
      \left(\frac{\hbar}{2}\right)
      \left[
        (\cos{\beta}-\frac{2}{\hbar}\lambda)\alpha_+ + \sin{\beta}e^{-i\alpha}\alpha_-
        \right] = (\cos{\beta}-\frac{2}{\hbar}\lambda)\alpha_+ + \sin{\beta}e^{-i\alpha}\alpha_-\\
      &0 =
      \left(\frac{\hbar}{2}\right)
      \left[
        \sin{\beta}e^{i\alpha}\alpha_+ - (\cos{\beta} + \frac{2}{\hbar}\lambda)\alpha_-
        \right] = \sin{\beta}e^{i\alpha}\alpha_+ - (\cos{\beta} + \frac{2}{\hbar}\lambda)\alpha_-\\
    \end{split}
  \end{equation}

  We put together our two sets of independent expressions for
  $\lambda$, $\alpha_+$, and $\alpha_-$, starting with the first one
  
  \begin{equation}
    \begin{split}
      0
      &= (\cos{\beta}-\frac{2}{\hbar}\lambda)\alpha_+ + \sin{\beta}e^{-i\alpha}\alpha_-\\
      &\Rightarrow
      0 = \cos{\beta}\alpha_+ - \frac{2}{\hbar}\lambda\alpha_+ + \sin{\beta}e^{-i\alpha}\alpha_-\\
      &\Rightarrow
      \frac{2}{\hbar}\lambda\alpha_+ = \cos{\beta}\alpha_+ + \sin{\beta}e^{-i\alpha}\alpha_-\\
      &\Rightarrow
      \lambda_1 = \frac{\hbar}{2}\frac{\cos{\beta}\alpha_+ + \sin{\beta}e^{-i\alpha}\alpha_-}{\alpha_+}\\
    \end{split}
  \end{equation}

  \begin{equation}
    \begin{split}
      0
      &= (\cos{\beta}-\frac{2}{\hbar}\lambda)\alpha_+ + \sin{\beta}e^{-i\alpha}\alpha_-\\
      &\Rightarrow
      0 = \frac{(\cos{\beta}-\frac{2}{\hbar}\lambda)\alpha_+}{\cos{\beta}-\frac{2}{\hbar}\lambda} + \frac{\sin{\beta}e^{-i\alpha}\alpha_-}{\cos{\beta}-\frac{2}{\hbar}\lambda}\\
      &\Rightarrow
      \alpha_+ = - \frac{\sin{\beta}e^{-i\alpha}}{\cos{\beta}-\frac{2}{\hbar}\lambda}\alpha_-\\
      &\Rightarrow
      \alpha_- = - \frac{\cos{\beta}-\frac{2}{\hbar}\lambda}{\sin{\beta}e^{-i\alpha}}\alpha_+
    \end{split}
  \end{equation}

  Moving on to our second independent equation and solving for those same 3 values
  
  \begin{equation}
    \begin{split}
      0
      &= \sin{\beta}e^{i\alpha}\alpha_+ - (\cos{\beta} + \frac{2}{\hbar}\lambda)\alpha_-\\
      &\Rightarrow
      0 = \sin{\beta}e^{i\alpha}\alpha_+ - \cos{\beta}\alpha_- - \frac{2}{\hbar}\lambda\alpha_-\\
      &\Rightarrow
      \frac{2}{\hbar}\lambda\alpha_- = \sin{\beta}e^{i\alpha}\alpha_+ - \cos{\beta}\alpha_-\\
      &\Rightarrow
      \lambda_2 = \frac{\hbar}{2}\frac{\sin{\beta}e^{i\alpha}\alpha_+ - \cos{\beta}\alpha_-}{\alpha_-}\\
    \end{split}
  \end{equation}
  
  \begin{equation}
    \begin{split}
      0
      &= \sin{\beta}e^{i\alpha}\alpha_+ - (\cos{\beta} + \frac{2}{\hbar}\lambda)\alpha_-\\
      &\Rightarrow
      0 = \frac{\sin{\beta}e^{i\alpha}\alpha_+}{\sin{\beta}e^{i\alpha}} - \frac{(\cos{\beta} + \frac{2}{\hbar}\lambda)\alpha_-}{\sin{\beta}e^{i\alpha}}\\
      &\Rightarrow
      \alpha_+ = \frac{\cos{\beta} + \frac{2}{\hbar}\lambda}{\sin{\beta}e^{i\alpha}}\alpha_-\\
      &\Rightarrow
      \alpha_- = \frac{\sin{\beta}e^{i\alpha}}{\cos{\beta} + \frac{2}{\hbar}\lambda}\alpha_+\\
    \end{split}
  \end{equation}

\noindent And when [t]his work is done -

\noindent Haha! - begins the fun.

\noindent From two equations

\noindent To one giant mess 

\noindent By way of Iliysk,

\noindent And Novorossiysk,

\noindent To Alexandrovsk to Akmolinsk

\noindent To Tomsk to Omsk

\noindent To Pinsk to Minsk

\noindent To here Eigenvalues run,

\noindent Yes, to here th'Eigenvalues run!

  \begin{equation}
    \begin{split}
      \alpha_+ &= - \frac{\sin{\beta}e^{-i\alpha}}{\cos{\beta}-\frac{2}{\hbar}\lambda}\alpha_-
      = - \frac{\sin{\beta}e^{-i\alpha}}{\cos{\beta}-\frac{2}{\hbar}\lambda}\frac{\sin{\beta}e^{i\alpha}}{\cos{\beta} + \frac{2}{\hbar}\lambda}\alpha_+\\
      &\Rightarrow
      1 = - \frac{\sin{\beta}e^{-i\alpha}}{\cos{\beta}-\frac{2}{\hbar}\lambda}\frac{\sin{\beta}e^{i\alpha}}{\cos{\beta} + \frac{2}{\hbar}\lambda}\\
      &\Rightarrow
      1 = - \frac{\cos{\beta}-\frac{2}{\hbar}\lambda}{\sin{\beta}e^{-i\alpha}}\frac{\cos{\beta} + \frac{2}{\hbar}\lambda}{\sin{\beta}e^{i\alpha}}\\
      &\Rightarrow
      1 = - \left(\cos{\beta}-\frac{2}{\hbar}\lambda\right)\left(\cos{\beta} + \frac{2}{\hbar}\lambda\right) \frac{1}{\sin^2{\beta}}\\
      &\Rightarrow
      1 = - \frac{\cos^2{\beta} + \frac{4}{\hbar^2}\lambda^2}{\sin^2{\beta}}\\
      &\Rightarrow
      \sin^2{\beta} = - \cos^2{\beta} + \frac{4}{\hbar^2}\lambda^2\\
      &\Rightarrow
      \sin^2{\beta} + \cos^2{\beta} = \frac{4}{\hbar^2}\lambda^2\\
      &\Rightarrow
      1 = \frac{4}{\hbar^2}\lambda^2\\
      &\Rightarrow
      \frac{\hbar^2}{4} = \lambda^2\\
      &\Rightarrow
      \lambda = \pm \frac{\hbar}{2}\qed\\
    \end{split}
  \end{equation}

Putting in our Eigenvalues, and solving for $\alpha_\pm$
  
  \begin{equation}
    \begin{split}
      &0 = (\cos{\beta}-\frac{2}{\hbar}\lambda)\alpha_+ + \sin{\beta}e^{-i\alpha}\alpha_-
      = (\cos{\beta}-1)\alpha_+ + \sin{\beta}e^{-i\alpha}\alpha_-\\
      &0 = \sin{\beta}e^{i\alpha}\alpha_+ - (\cos{\beta} + \frac{2}{\hbar}\lambda)\alpha_-
      = \sin{\beta}e^{i\alpha}\alpha_+ - (\cos{\beta} + 1)\alpha_-\\
    \end{split}
  \end{equation}

  \begin{equation}
    \begin{split}
      \alpha_+ &= -\frac{\sin{\beta}e^{-i\alpha}\alpha_-}{\cos{\beta}-1}\\
      \alpha_- &= \frac{\sin{\beta}e^{i\alpha}\alpha_+}{\cos{\beta} + 1}\\
    \end{split}
  \end{equation}

  When we start setting up to build our Eigenvectors, we're going to
  need these half-angle formulas:
  
  \begin{equation}
    1 + \cos{\beta} = 2 \cos^2{\frac{\beta}{2}}
  \end{equation}

  \begin{equation}
    1 - \cos{\beta} = 2 \sin^2{\frac{\beta}{2}}
  \end{equation}

  My humor largely exhausted, and soul much relieved to have done so,
  I present unto you, patient (and likely mortified) reader, the final
  derivation of our eigenvector:
  
  \begin{equation}
    \begin{split}
      \frac{\alpha_+}{\alpha_-}
      &= -\frac{\sin{\beta}e^{-i\alpha}\alpha_-}{\cos{\beta}-1}
      \frac{\cos{\beta} + 1}{\sin{\beta}e^{i\alpha}\alpha_+}\\
      &= -\frac{\alpha_-}{\alpha_+}\frac{\sin{\beta}e^{-i\alpha}}{\sin{\beta}e^{i\alpha}}
      \frac{\cos{\beta} + 1}{\cos{\beta} - 1}\\
      &= -\frac{\alpha_-}{\alpha_+}\frac{e^{-i\alpha}}{e^{i\alpha}}
      \frac{\cos{\beta} + 1}{\cos{\beta} - 1}\\
      &= -\frac{\alpha_-}{\alpha_+}e^{-2i\alpha}
      \frac{\cos{\beta} + 1}{\cos{\beta} - 1}\\
      \frac{\alpha_+^2}{\alpha_-^2}
      &= -e^{-2i\alpha}\frac{\cos{\beta} + 1}{\cos{\beta} - 1}\\
      &= e^{-2i\alpha}\frac{1 + \cos{\beta}}{1 - \cos{\beta}}\\
      &= e^{-2i\alpha}\frac{2 \cos^2{\frac{\beta}{2}}}{2 \sin^2{\frac{\beta}{2}}}\\
      &= e^{-2i\alpha}\frac{\cos^2{\frac{\beta}{2}}}{\sin^2{\frac{\beta}{2}}}\\
      \frac{\alpha_+}{\alpha_-}
      &= e^{-i\alpha}\frac{\cos{\frac{\beta}{2}}}{\sin{\frac{\beta}{2}}}\\
    \end{split}
  \end{equation}

  \begin{equation}
    \ket{\alpha_+} = \cos{\frac{\beta}{2}} \ket + + e^{-i\alpha}\sin{\frac{\beta}{2}} \ket - \qed
  \end{equation}

  The appropriate next step is Laguvulin 16, or maybe Bowmore 18, or,
  you know, some acetone or some stricnine or something.

\end{answer}
